\chapter{Descripción del Proyecto}
\label{ch:overview}
% ==============================================================

\section{Contexto y motivación}

La playa de Guardamar del Segura es un espacio de alto valor ecológico y turístico sujeto a
procesos de erosión y acreción que requieren seguimiento sistemático. La medición manual de la
línea de costa mediante levantamientos GNSS es costosa, puntual y no escala para cubrir todo el
frente litoral con frecuencia horaria.

El proyecto \textbf{cv-lit} surge para automatizar esta medición aprovechando la infraestructura
de cámaras fijas ya desplegada por el servicio de vigilancia de la playa, convirtiendo cada
fotografía en un dato georreferenciado de la posición de la interfaz arena--mar.

\section{Partes implicadas}

\begin{table}[H]
\centering
\caption{Actores del proyecto}
\label{tab:actors}
\begin{tabular}{L{4cm} L{9cm}}
\toprule
\textbf{Rol} & \textbf{Entidad} \\
\midrule
Cliente & Ayuntamiento de Guardamar del Segura \\
Datos GNSS (GCPs) & Instituto de Ecología Litoral (IEL) \\
Desarrollo & 1 ingeniero a tiempo parcial (Tech4D Lab, UA) \\
Supervisión & Universidad de Alicante \\
Duración & 6 meses (enero -- junio 2026) \\
\bottomrule
\end{tabular}
\end{table}

\section{Objetivos del sistema}

El producto final es un \textbf{GeoJSON proyectado en EPSG:25830} con los siguientes atributos
obligatorios por cada detección:

\begin{table}[H]
\centering
\caption{Atributos del GeoJSON de salida}
\label{tab:geojson-fields}
\begin{tabular}{L{3.5cm} L{3.5cm} L{5.5cm}}
\toprule
\textbf{Campo} & \textbf{Tipo} & \textbf{Descripción} \\
\midrule
\texttt{ID\_Camara}    & Integer (1--6) & Identificador de cámara origen \\
\texttt{Timestamp}     & ISO 8601 UTC   & Fecha y hora de captura \\
\texttt{Confianza\_IA} & Float (0--1)   & Score de confianza del modelo SAM \\
\texttt{Area\_Seca\_m2}& Float          & Superficie de playa seca (m²) \\
\bottomrule
\end{tabular}
\end{table}

\section{Restricciones no negociables}

\begin{itemize}[leftmargin=1.5em]
  \item Toda implementación \textbf{debe} producir GeoJSON en \textbf{EPSG:25830}.
  \item Fase horaria prioritaria: \textbf{12:00h solar local}.
  \item Modelo de segmentación: \textbf{SAM} (Segment Anything Model, Meta AI).
  \item Cada cámara tiene su propio perfil de calibración \textit{independiente}.
  \item Criterio de aceptación de precisión: \textbf{RMSE~$<$~1,5~m} en condiciones estándar.
\end{itemize}

\section{Entregables}

\begin{enumerate}[leftmargin=1.5em]
  \item Aplicación ejecutable configurada y documentada (esta memoria).
  \item Perfiles de calibración por cámara (6~ficheros \texttt{.npy} / \texttt{.json}).
  \item Interfaz web de operación y validación por lotes.
  \item Informe de validación con métricas RMSE y MAE.
\end{enumerate}

% ==============================================================
